\documentclass[11pt]{article}
\usepackage[margin=1in]{geometry}
\usepackage{amsmath,amssymb,booktabs,array}
\usepackage[hidelinks]{hyperref}
\usepackage{enumitem}
\title{EmbodiK Explicit Task Stacks: Formulation, API, and Literature Context}
\author{EmbodiK developer note}
\date{September 2026}

\begin{document}
\maketitle

\begin{abstract}
This note describes the named task-stack API being added to EmbodiK and how it
relates to the existing eSNS-style hierarchical solver. The implementation
does not introduce a new inverse or hard-constraint algorithm. It assembles
registered objectives by named priority level and routes them through the
existing constrained solver and singularity robust inverse (SRI). The optional
\texttt{LEXICOGRAPHIC\_LEAST\_SQUARES} stack policy is a convenience mapping to
the existing \texttt{MIN\_ERROR} objective mode at every level. Global hard
constraints remain separate from task objectives.
\end{abstract}

\section{Problem statement}

Let $v\in\mathbb{R}^{n_v}$ be the commanded generalized velocity. The explicit
stack contains $K$ levels, ordered from highest to lowest priority. After
assembling same-level registered tasks, level $i$ is represented by
\begin{equation}
  J_i v \simeq b_i,
  \qquad J_i\in\mathbb{R}^{m_i\times n_v},\quad
  b_i\in\mathbb{R}^{m_i}.
  \label{eq:level}
\end{equation}
All member tasks at a level are peers: their Jacobians and target velocities
are row-stacked to form $(J_i,b_i)$. Task names are canonicalized for
deterministic assembly. The level sequence, not the integer priority stored on
each task, determines hierarchy while explicit stack configuration is active.

Hard constraints are represented independently as a feasible set
\begin{equation}
  \mathcal{F}=\{v\mid \ell_h\le C_hv\le u_h\}.
  \label{eq:hard}
\end{equation}
The global hard rows can include joint velocity/position limits, collision,
CoM, relative-pose, user linear bounds, and active contact conditions. They
are passed through the existing constrained solve path. They are not appended
as task levels and are not converted into residual costs. This is a deliberate
API boundary: every priority level is an objective, while hard feasibility is
shared across the whole hierarchy.

\section{Two objective policies already supported}

\subsection{Direction-preserving SCALE}
The SNS-style scaling policy seeks the largest feasible scalar
$s_i\in[0,1]$ for the requested target direction, subject to the current
constraints and higher-level solution:
\begin{equation}
  \max_{v,s_i}\;s_i
  \quad\text{such that}\quad J_i v=s_i b_i,\quad
  v\in\mathcal{F}_i,\quad 0\le s_i\le 1.
  \label{eq:scale}
\end{equation}
Here $\mathcal{F}_i$ also reflects the hierarchy already established by
levels $1,\ldots,i-1$. When the complete target is infeasible, reducing $s_i$
retains its direction while sacrificing magnitude.

\subsection{Minimum residual MIN\_ERROR}
The existing minimum-error policy instead minimizes the level residual:
\begin{equation}
  r_i^*=\min_{v\in\mathcal{F}_{i-1}}
       \|J_i v-b_i\|_2,
  \qquad
  \mathcal{F}_i=\{v\in\mathcal{F}_{i-1}:
       \|J_i v-b_i\|_2\le r_i^*+\epsilon_i\}.
  \label{eq:minerr}
\end{equation}
Equation~\eqref{eq:minerr} is the ideal lexicographic least-squares statement:
each level minimizes its error without worsening the achieved higher-level
solutions. In the implementation, SRI damping, rank thresholds, active
constraints, and finite precision determine the numerical approximation and
the effective preservation tolerance. Thus the API promises the existing
hierarchical solver's behavior and status reporting; it does not claim exact
arithmetic optimality for every ill-conditioned case.

With an unreachable target, $r_i^*>0$ can be valid. A lower level may use
remaining directions, but should not trade away a higher-level achieved
solution. Global hard infeasibility is a solver infeasibility and is distinct
from a nonzero task residual.

The two policies differ in their local objective. SCALE preserves the
requested target direction as much as feasibility permits. MIN\_ERROR can use
reachable components even when the full target direction is infeasible. Both
are objective policies in the existing hierarchical solver; MIN\_ERROR is not
a new backend added by task stacks.

\section{What the stack API adds}

The current C++ and Python surface is:
\begin{verbatim}
TaskLevelSpec(name, task_names, solve_mode=SCALE,
              allow_min_error_fallback=false)
TaskStackConfig(levels, backend=SNS)
solver.configure_task_stack(config)
solver.task_stack_config
solver.has_explicit_task_stack()
solver.clear_task_stack()
\end{verbatim}

The C++ types are \texttt{embodik::TaskLevelSpec} and
\texttt{embodik::TaskStackConfig}; Python exposes them as
\texttt{TaskLevelSpec} and \texttt{TaskStackConfig}. The selector applies to
velocity solves and registered-task position steps. Standalone position
solves build their own objective sequence.

The default \texttt{SNS} policy retains each level's configured mode. The
explicit least-squares policy maps each level to the existing \texttt{MIN\_ERROR}
mode and uses the same multi-objective solve route, global constraint rows,
active-set machinery, and SRI as the default route. It is a named convenience
preset, not a separate dense QP solver. \texttt{SCALE\_ELASTIC}
and per-level SNS fallback flags are rejected under that preset rather than
silently discarded.

The stack additionally defines observable behavior for membership and
configuration: levels are explicitly ordered; same-level tasks form one
objective; tasks omitted from the explicit stack are omitted from that
registered-task solve; member task priorities and policies are not mutated;
and stale/missing members are reported as invalid input.

\subsection{Runtime push/pop for agentic applications}

The current API replaces the complete value with
\texttt{configure\_task\_stack}; it does not yet provide scoped
\texttt{push}/\texttt{pop} operations. A safe extension is an immutable
snapshot plus validated patches, for example:
\begin{verbatim}
patch = TaskStackPatch(base_revision=solver.task_stack_revision,
                       source="planner-session-42",
                       push_levels=[level_spec])
receipt = solver.apply_task_stack_patch(patch)
solver.remove_task_stack_patch(receipt.patch_id)
\end{verbatim}
The shorthand \texttt{push\_level(spec) -> token} and \texttt{pop(token)}
should be wrappers around that patch mechanism. A blind global \texttt{pop()}
is ambiguous when planner, operator, and recovery sources can all edit a
stack.

For this interface to work well with agentic APIs, it should have the
following contract:
\begin{enumerate}[leftmargin=*,nosep]
  \item Each patch names a base revision; accepted edits produce a monotonic
        revision and stale edits are rejected instead of overwriting a newer
        intent.
  \item Validation is transactional: either the full candidate configuration
        is valid and published, or the active snapshot is unchanged.
  \item Each patch has an opaque ID, source, and optional idempotency key.
        Removing a patch reverses only that patch's contribution.
  \item A solve reads one stable immutable snapshot. Updates become visible
        between control ticks, never halfway through a solve.
  \item Optional leases expire temporary intents to a deterministic fallback
        snapshot if the source disconnects or stops renewing.
  \item Agents submit bounded task intents or registered task names. They do
        not submit arbitrary Jacobians, constraints, executable code, or
        changes to global hard bounds. A trusted policy layer decides whether
        an intent may be committed.
\end{enumerate}
The update receipt should make the state observable: accepted/rejected,
active revision, patch ID, and a validation or conflict reason. The solver
should not silently accept stale or partially valid edits. These are
architectural recommendations, not APIs implemented in this branch.

\subsection{A model-neutral whole-body controller contract}

The task-stack API configures the solver. VLA, WAM, and general-purpose agent
callers should usually use a higher-level typed controller API instead of
building Jacobians or manipulating priority directly. A useful request record
has the form
\begin{equation}
  g=(\text{capability},\text{target},\text{frame},\text{tolerances},
      \text{contact/approach},\text{time window},\text{source},\text{expiry}),
\end{equation}
where the capability is an allowed operation such as reach, look-at, hold, or
bimanual grasp, and the target is expressed in a named frame with explicit
units. A trusted mapper resolves the capability to registered tasks and an
approved stack template. The agent does not choose the global hard bounds.

The controller contract should then provide:
\begin{enumerate}[leftmargin=*,nosep]
  \item \textbf{Preflight} that checks target/frame validity, current robot
        capability, and constraints, returning structured reasons and
        relevant limits rather than only a Boolean.
  \item \textbf{Execution handles} with \texttt{status}, \texttt{cancel},
        and \texttt{revise}; action proposals have a bounded time horizon and
        are revalidated as state changes.
  \item \textbf{Feedback events} containing state timestamps, active stack
        revision, task residuals, binding hard constraints, progress, and a
        terminal reason such as reached, blocked, infeasible, stale target,
        or cancelled.
  \item \textbf{Structured traces} joining intent, revision, solver result,
        executed interval, and outcome verification. Model/prompt identity
        belongs in the orchestration layer and should not be hard-coded into
        the numerical solver.
\end{enumerate}

Keep the agent loop asynchronous and lower-rate than the deterministic
whole-body control loop. VLA action chunks and WAM-predicted trajectories are
bounded proposals: preflight them, execute a short segment, observe the
result, then replan. This is compatible with different model providers and
versions (including models referred to as Astra or Fable) without coupling
the controller ABI to any one model's tool schema. Permission to propose an
intent and permission to commit physical motion should be separate policy
decisions. Recent VLA-agent work makes prerequisite checks, bounded execution,
post-action verification, and structured recovery traces explicit
\cite{embodiedskills2026}; WAM work emphasizes compatible video/action
interfaces when composing independently trained components
\cite{openwam2026}. A published Astra/Fable robot-control trial exposes
camera/proprioceptive observations and end-effector pose/gripper actions via
an agent tool interface, though its reported setup has experimental limits
\cite{robottoolapi2026}. These examples motivate an interoperable controller
contract, not model-specific solver APIs.

\section{Results and caller expectations}

For an explicit prioritized solve, velocity and position-step results expose
per-level diagnostics when the backend returns a one-to-one result:
\begin{itemize}[nosep]
  \item level name and canonical member task names;
  \item configured and effective solve modes;
  \item task scale, residual norm, requested target norm, and normalized
        residual;
  \item hierarchy backend and accepted solve path;
  \item prioritized status before any constrained weighted fallback.
\end{itemize}
For MIN\_ERROR, scale is reported as $-1$. Residuals describe the combined
level, not separate member tasks. If weighted recovery replaces a failed
prioritized result, the result path says \texttt{WEIGHTED\_FALLBACK},
preservation is reported false, and per-level diagnostics are empty.

Callers should use both the top-level status and the per-level residuals. A
successful solve can have nonzero error on an unreachable target. Conversely,
global hard infeasibility must not be interpreted as ordinary task degradation.

\section{Relation to the literature}

The stack-of-tasks (SoT) idea is a representation of ordered objectives, not a
single unique numerical algorithm. Classical null-space methods, SNS/eSNS,
hierarchical quadratic programming (HQP), and sequential hierarchical
least-squares methods differ in how they enforce hierarchy, handle
inequalities, and trade numerical accuracy against computation.

\begin{center}
\small
\begin{tabular}{>{\raggedright\arraybackslash}p{0.23\linewidth}
                >{\raggedright\arraybackslash}p{0.34\linewidth}
                >{\raggedright\arraybackslash}p{0.34\linewidth}}
\toprule
\textbf{Approach} & \textbf{Hierarchy and constraints} &
\textbf{Relation to this API} \\
\midrule
SNS/eSNS & Saturation-based redundancy resolution; eSNS generalizes the
formulation to equality and inequality constraints at multiple levels. &
Closest algorithmic family. EmbodiK reuses its existing constrained
multi-objective solver and SRI path. \\
HQP & A sequence of constrained quadratic programs can place inequality
constraints at priority levels. & More general constraint hierarchy than
this API, which keeps hard constraints global. \\
Sequential hierarchical least squares & Directly solves ordered least-squares
objectives; recent methods study speed, numerical stability, and optimality. &
Conceptually describes the MIN\_ERROR hierarchy; no independent solver is
added here. \\
Explicit stack surface & Names, memberships, order, modes, result diagnostics,
and scoped configuration. & Configuration/API layer around the current
solver, not a new algorithm family. \\
\bottomrule
\end{tabular}
\end{center}

Flacco et al.'s SNS handles hard joint motion bounds and prioritized tasks by
saturating constraints in the task null space; it integrates task scaling
when the requested task is infeasible~\cite{flacco2015}. Fiore et al.'s eSNS
generalizes the SNS formulation to arbitrary equality and inequality
constraints at multiple priority levels, control-effort weights, and residual
redundancy inputs, with basic, fast, and optimality-oriented
variants~\cite{fiore2022}. Escande, Mansard, and Wieber formulate strict
hierarchies of equality and inequality least-squares problems using HQP
machinery~\cite{escande2014}. Recent sequential hierarchical least-squares
work studies solver efficiency, infeasible constraints, numerical stability,
and optimality for robot control and planning~\cite{chefchaouni2024}.

This literature comparison yields an important boundary for this API:
eSNS/HQP formulations can place inequality constraints at priority levels,
whereas EmbodiK's explicit stack keeps its hard constraints global and
independent of task order. This is consistent with the current architecture
and the desired separation between hard constraints and prioritized tasks,
but it is a narrower specification than full hierarchical constraint
programming.

\section{Potential gaps to resolve}

The literature suggests several areas to make explicit before describing the
API as a general-purpose SoT solver:
\begin{enumerate}[leftmargin=*,nosep]
  \item \textbf{Peer-task scaling.} Equation~\eqref{eq:minerr} uses the
        assembled row residual. Position, rotation, momentum, and posture
        errors have different units and magnitudes. The current stack surface
        does not define a general per-task normalization/weight policy for
        every objective family. Specify weights or a normalization contract
        before recommending heterogeneous tasks as peers.
  \item \textbf{Constraint hierarchy.} There is no per-level inequality
        objective/constraint API. Keep safety and hard bounds globally active;
        only add constraint priority levels for a concrete use case with
        explicit feasibility and relaxation semantics.
  \item \textbf{Numerical strictness.} SRI damping and rank/feasibility
        thresholds make practical preservation approximate near singularities.
        Document the existing tolerances, and distinguish nominal hierarchy
        semantics from finite-precision guarantees.
  \item \textbf{Stack transitions.} Changing order or membership between
        control ticks can change the command discontinuously. Applications
        that switch stacks may need continuity logic or state-transition
        smoothing.
  \item \textbf{Runtime evidence.} The named stack inherits existing solver
        costs and behavior. It needs no independent real-time claim; any
        performance statement should be measured against the current solver
        on representative problem sizes.
  \item \textbf{Validation coverage.} Cover same-level unit scaling,
        conflicting/unreachable levels, redundant and rank-deficient rows,
        hard-bound infeasibility, near-singular Jacobians, and transitions
        between legacy and explicit stacks.
\end{enumerate}

These are API clarity and validation gaps, not evidence that a new optimizer
is missing. The implementation choice is to reuse the established solver and
make its existing objective modes easier to configure and inspect.

\begin{thebibliography}{9}
\bibitem{flacco2015}
F. Flacco, A. De Luca, and O. Khatib, ``Control of Redundant Robots Under
Hard Joint Constraints: Saturation in the Null Space,'' \emph{IEEE
Transactions on Robotics}, vol. 31, no. 3, pp. 637--654, 2015.
\href{https://doi.org/10.1109/TRO.2015.2418582}{doi:10.1109/TRO.2015.2418582}.
\href{https://www.diag.uniroma1.it/~labrob/pub/papers/TRO_SNS_Jun2015.pdf}{Author PDF}.

\bibitem{fiore2022}
M. D. Fiore, G. Meli, A. Ziese, B. Siciliano, and C. Natale, ``A General
Framework for Hierarchical Redundancy Resolution Under Arbitrary Constraints,''
\emph{IEEE Transactions on Robotics}, 2023.
\href{https://arxiv.org/abs/2204.03974}{arXiv:2204.03974}.

\bibitem{escande2014}
A. Escande, N. Mansard, and P.-B. Wieber, ``Hierarchical Quadratic
Programming: Fast Online Humanoid-Robot Motion Generation,'' \emph{The
International Journal of Robotics Research}, vol. 33, no. 7, pp. 1006--1028,
2014. \href{https://doi.org/10.1177/0278364914521306}{doi:10.1177/0278364914521306}.
\href{https://hal.science/hal-00751924}{HAL record}.

\bibitem{chefchaouni2024}
S. Chefchaouni, M. Benallegue, A. Escande, and P.-B. Wieber, ``Efficient
Lexicographic Optimization for Prioritized Robot Control and Planning,'' 2024.
\href{https://arxiv.org/abs/2403.09160}{arXiv:2403.09160}.

\bibitem{openvla2025}
M. J. Kim et al., ``OpenVLA: An Open-Source Vision-Language-Action Model,''
\emph{Proceedings of the 8th Conference on Robot Learning}, PMLR vol. 270,
2025. \href{https://proceedings.mlr.press/v270/kim25c.html}{PMLR paper}.

\bibitem{openwam2026}
OpenWAM Team, ``OpenWAM: An Open Framework for Composable World-Action
Models,'' Stanford University project and technical overview, 2026.
\href{https://openwam.stanford.edu/}{Project overview}.

\bibitem{embodiedskills2026}
``EmbodiedSkills: A Unified Framework for Orchestrating, Training, and
Deploying VLA Agents,'' arXiv:2609.01281, 2026.
\href{https://arxiv.org/abs/2609.01281}{arXiv paper}.

\bibitem{robottoolapi2026}
Robocurve, ``GPT-6 Astra on robot arms,'' technical setup and trial report,
2026. Its documented agent interface uses camera/proprioceptive observations
and absolute end-effector pose plus gripper commands; the page also reports
trial and setup limitations.
\href{https://openai.robocurve.org/gpt-6-astra/}{Technical report}.
\end{thebibliography}

\end{document}
